\section{Results}
\label{sec:results}

% Every number in this section was read from the cited frozen artifact during the
% 2026-08-07 audit. Numbers from superseded runs are NOT carried forward; see
% manuscript/audit/DRAFT_STALENESS.md.

\subsection{Atlas construction and data roles}
\label{sec:res-atlas}

% RESULT SOURCE:
% results/motif_pipeline/core3_atlas_v3/data/v3_region_manifest.csv
% results/motif_pipeline/core3_atlas_v3/data/V3_PREDECLARATION.json
Three structural families are represented: CsCl, Th$_3$P$_4$ and amorphous. Two training
atlases were constructed on an identical condition-group split, differing only in training
content. Atlas~A contains only expert-supported labels; atlas~B additionally admits
machine-generated provisional labels, confined to training. Calibration and validation
membership are identical between the two and entirely expert-labelled, so any difference in
outcome is attributable to training content alone.

\begin{table}[htbp]
  \centering\small
  \caption{Data roles. Calibration and validation are shared between the two atlases.}
  \label{tab:splits}
  \begin{tabular}{lrrrr}
    \toprule
    role & atlas A regions & atlas B regions & condition groups & label source \\
    \midrule
    training    & 126 & 260 & 68 / 70 & expert (A); expert + provisional (B) \\
    calibration & \multicolumn{2}{c}{70 (shared)} & 19 & expert only \\
    validation  & \multicolumn{2}{c}{50 (shared)} & 14 & expert only \\
    untouched test & \multicolumn{2}{c}{sealed} & 6 & not opened \\
    \bottomrule
  \end{tabular}
\end{table}

Every trajectory entering the atlas was re-processed through the frozen gas and
reliable-centre screens, with the gates asserted rather than assumed. On the subset that
overlaps a previously built descriptor set, the re-derived descriptors are bit-identical,
confirming that the representation is reproducible across machines.
% RESULT SOURCE:
% results/motif_pipeline/core3_atlas_v3/data/V3_SHARDS_COMPLETE.json

\subsection{Unseen structural families are detected}
\label{sec:res-lofo}

% RESULT SOURCE:
% results/motif_pipeline/core3_atlas_v3/data/v3_m1_lofo.csv
Under the strict protocol of \cref{sec:lofo}, in which an entire family is removed and all
statistical quantities re-fitted, both atlases separate the held-out family perfectly in
ranking: AUROC and true-positive rate at the calibrated threshold are $1.000$ for all three
families. The discriminating quantity is the false-positive rate on known domains at the
same threshold (\cref{tab:lofo}).

\begin{table}[htbp]
  \centering\small
  \caption{Strict leave-one-family-out. AUROC and TPR at $\tau$ are $1.000$ in every row;
  the false-positive rate on retained families is what differs.}
  \label{tab:lofo}
  \begin{tabular}{lrrrr}
    \toprule
    held-out family & positives & negatives & FPR at $\tau$ (A) & FPR at $\tau$ (B) \\
    \midrule
    CsCl        & 31 & 37 & 0.027 & 0.027 \\
    Th$_3$P$_4$ & 27 & 39 & 0.103 & 0.077 \\
    amorphous   & 62 & 24 & 0.083 & 0.292 \\
    \bottomrule
  \end{tabular}
\end{table}

That a fixed descriptor, with no exposure to the held-out family at any stage, ranks unseen
domains above all retained ones is the central detection result. We note that perfect
ranking on evaluation sets of $61$--$86$ domains does not establish perfect ranking in
general, and the false-positive column shows the threshold is doing real work even where the
ranking is clean.

\subsection{Provisional training labels do not help}
\label{sec:res-ab}

% RESULT SOURCE:
% results/motif_pipeline/core3_atlas_v3/data/v3_m1_lofo.csv
% results/motif_pipeline/core3_atlas_v3/data/v3_m2_false_positive.csv
% results/motif_pipeline/core3_atlas_v3/data/v3_m4_threshold_stability.csv
Atlas~B carries $134$ additional training regions. Against a predeclared set of five
metrics it produces no measurable improvement and one clear regression.

On leave-one-family-out, B is better on the Th$_3$P$_4$ holdout by a single domain
($0.077$ versus $0.103$) and substantially worse on amorphous ($0.292$ versus $0.083$).
The known false-positive rate is indistinguishable ($0.22$ versus $0.24$, a difference of
one domain in fifty). Clean-control retention and centre-set sensitivity are identical.
Threshold stability separates them decisively: resampling calibration groups gives a
relative spread in $\tau$ of $0.031$ for A against $0.098$ for B, and a validation decision
flip fraction of $0.0063$ against $0.0234$. Since both atlases calibrate on the same $19$
groups, this difference originates entirely in training content.

The failure concentrates where the provisional labels concentrate. Amorphous accounts for
$109$ of B's $260$ training regions against $52$ of A's $126$, and amorphous is also the
class on which expert adjudication most often overturned the machine labels.
\TODO{VERIFY: quantify the adjudication overturn rate by class from the adjudication record
rather than describing it qualitatively.}

\subsection{The acceptance threshold, not the score, is the fragile component}
\label{sec:res-threshold}

% RESULT SOURCE:
% results/motif_pipeline/core3_atlas_v3/data/v3_threshold_diagnostic.csv
% results/motif_pipeline/core3_atlas_v3/data/v3_threshold_rule_options.csv
The predeclared threshold rule---the $95$th percentile of per-group mean calibration
scores---yields a known false-positive rate of $0.24$ (A) and $0.22$ (B) on in-distribution
validation domains. This is high, and the diagnostic identifies why: the statistic that sets
$\tau$ is a group mean, while the statistic it is compared against is an individual domain
score. Averaging within a group removes the within-group spread, so the threshold is
narrower than the distribution it gates.

The effect is quantified by costing alternative rules on the same data. Taking instead the
$99$th percentile of individual calibration scores reduces the false-positive rate to
$0.040$ (A) and $0.060$ (B). No rule is adopted here; the predeclared rule is retained for
the comparison of \cref{sec:res-ab}, and the choice is left open.

\subsection{A control panel that no longer discriminates}
\label{sec:res-controls}

% RESULT SOURCE:
% results/motif_pipeline/core3_atlas_v3/data/v3_m3_clean_controls.csv
% results/motif_pipeline/core3_atlas_v3/data/v3_threshold_diagnostic.csv
A panel of $24$ previously constructed control domains---eight built to be clean
realisations and sixteen deliberately perturbed---was re-scored. Both atlases reject all
sixteen perturbed controls and also reject all eight clean ones.

This is not a threshold artefact. The clean controls have a median score of $9.14$ (A) and
$8.15$ (B), while the \emph{maximum} individual calibration score is $7.27$ and $6.75$
respectively: the clean controls lie beyond the entire calibration range, so no
calibration-derived threshold can accept them. Of six alternative rules costed, five give
zero of eight in both atlases.

The mechanism is visible in the fitted shrinkage. In the present configuration the
Ledoit--Wolf coefficient is $\approx 0.006$, so the covariance is essentially empirical and
the acceptance region is correspondingly tight. An earlier configuration that fitted the
same estimator in a regime where the covariance was ill-conditioned produced a much broader
acceptance region, and accepted seven of the eight clean controls. We interpret the earlier
acceptance as a consequence of a degenerate covariance rather than as recognition, noting
that the same earlier configuration rejected only six of sixteen perturbed controls.

Separately, re-scoring each validation domain under alternative centre sets shows the
detector is closely tied to the reliable-centre convention: $0.76$ (A) and $0.78$ (B) of
validation domains are accepted using the frozen centres, and none are accepted under any
alternative centre set tested. Because the control panel was constructed by a different
segmentation procedure, it carries a different centre set by construction, and we therefore
regard the clean-control panel as measuring convention mismatch rather than detector
quality in its present form.
\TODO{Rebuild the clean-control panel on the current centre convention before it is used as
an acceptance criterion.}

\subsection{Physical identification of a detected candidate}
\label{sec:res-matcher}

% RESULT SOURCE:
% results/statistical_ood_benchmark_v2/reference_conditioned_matcher_v1/steven_match_results.csv
An external candidate that the statistical detectors reject is passed to the reference
matcher. It is matched to Ni$_3$Sn at $D = 0.347$ with a runner-up margin of $1.86$, inside
the acceptance cutoff of $0.40$. A second candidate is matched to Cu$_3$Au at $D = 0.373$
with margin $1.93$. Both are structures absent from the three-family training atlas, so
detection and identification are consistent: the detector reports "not in the atlas" and the
matcher supplies an identity from an independent library.

\subsection{Abstention when the correct reference is removed}
\label{sec:res-abstain}

% RESULT SOURCE:
% results/statistical_ood_benchmark_v2/reference_conditioned_matcher_v1/cscl_holdout_results.csv
% results/statistical_ood_benchmark_v2/reference_conditioned_matcher_v1/amorphous_negative_control.csv
% results/statistical_ood_benchmark_v2/reference_conditioned_matcher_v1/inter_family_distance_matrix.csv
With the full library, $23$ independent CsCl domains are matched to CsCl at
$D = 0.16$--$0.23$. With the CsCl reference deleted, all $23$ abstain: their nearest
surviving reference is zangenite at $D = 0.60$--$0.71$. That range lies inside the
distribution of distances between two \emph{different} library anchors
($0.56$--$1.38$, median $0.995$), so the nearest survivor is no closer than an unrelated
structure would be, and the abstention is principled rather than threshold-dependent. A
negative control of eight amorphous domains abstains at $D = 0.40$--$0.43$.

The narrowest boundary in the benchmark is between the accepted candidates
($D \le 0.373$) and the amorphous control ($D \ge 0.404$), a gap of $0.03$. This should be
regarded as the limiting margin of the present matcher configuration.

\paragraph{Caveat on the zangenite anchor.} The abstention result rests on the zangenite
reference, and the reference-construction limitation of \cref{sec:refbuild} changes that
anchor materially. The Ni$_3$Sn anchor underlying \cref{sec:res-matcher} is unchanged, so
the identity assignment is unaffected, but the runner-up for the first candidate is
zangenite and therefore its margin is conditional.
\TODO{Re-derive the reference-removal arm with the anisotropically relaxed zangenite anchor,
including a regenerated thermal envelope.}

\subsection{Incremental atlas expansion}
\label{sec:res-expand}

The workflow is designed so that a newly identified family can be incorporated by adding
reference structures and re-fitting statistical quantities, without retraining a learned
representation. A direct test would proceed as follows: construct an atlas excluding family
$X$; verify that an independent sample of $X$ is detected as novel; verify that the matcher
identifies it; add a small number of $X$ references; re-fit the statistical atlas; and
confirm both that a second independent sample of $X$ is then recognised as known and that
the false-positive rate on the original families is unchanged.

No such sequential experiment exists in the present repository, and the expansion stage is
therefore \textbf{not demonstrated}.
\TODO{RUN SEQUENTIAL DETECT-MATCH-ADD-RECOGNIZE EXPERIMENT}
